\documentclass[handout]{beamer} 
% change the alerted colour to blue
\setbeamercolor{alerted text}{fg=blue}

% theme split
\usepackage{beamerthemesplit}

% theme shadow
\usepackage{beamerthemeshadow}

% logo
\logo{\includegraphics[height=1cm]{iitblogo.pdf}}

\begin{document}
% sf family, bold font
\sffamily \bfseries

\title
[Presentation using \LaTeX\ and Beamer
\hspace{0.5cm}
\insertframenumber/\inserttotalframenumber]
{Spoken Tutorial: \\ Presentation using \LaTeX\ and
  Beamer}
\author
[Kannan Moudgalya]
{Kannan M. Moudgalya \\
  Indian Institute of Technology, Bombay \\
  kannan@iitb.ac.in}
\date{IIT Bombay \\
  3 November 2009}

\begin{frame}
   \titlepage
\end{frame}

\begin{frame}
\frametitle{Outline}
\begin{itemize}
\item Title page, author name, colour, logo, etc.
\item Minimal animation
\item Two columns
\item Figures and Tables
\item Equations
\item Verbatim
\item References for further reading
\end{itemize}
\end{frame}

\begin{frame}
\frametitle{Other Spoken Tutorials on  \LaTeX}
\begin{itemize}
\item The following spoken tutorials
  on \LaTeX\ are available at
  {\color{magenta}http://moudgalya.org}:
\begin{itemize}
\item What is compilation?
\item Letter writing
\item Report writing
\item Mathematical typesetting
\item Equations
\item Tables and figures
\item How to create bibliography?
\item Inside story of bibliography
\end{itemize}
\item Installation and use of \LaTeX\ on windows
  is explained at
{\color{magenta} \small \tt http://www.che.iitb.ac.in/faculty/km/spoken/LaTeX-Win.wmv}
\item If you are not comfortable in using \LaTeX,
  you may consider going through these spoken
  tutorials first
\item More permanent links for all of these should
  be available in the future at {\color{magenta}
    \tt http://fossee.in}
\end{itemize}
\end{frame}

% equations - from spoken tutorial on equations
\begin{frame}
\frametitle{Equations:  from Spoken Tutorial on
  Equations} 
Let us start with the model of an inverted
pendulum:
\begin{align*}
\frac d{dt} 
\begin{bmatrix} x_1 \\ x_2 \\ x_3 \\
  x_4 \end{bmatrix} & =
\begin{bmatrix} 
0 & 0 & 1 & 0 \\
0 & 0 & 0 & 1 \\
0 & -\gamma & 0 & 0 \\
0 & \alpha & 0 & 0
\end{bmatrix}
\begin{bmatrix} x_1 \\ x_2 \\ x_3 \\
  x_4 \end{bmatrix} +
\begin{bmatrix} 0 \\ 0 \\ -\delta \\
  -\beta \end{bmatrix} \Delta\mu
%
\intertext{Proportional, integral, derivative
  controller is most popular in industry.  It has
  three tuning parameters: $K$, $\tau_i$ and
  $\tau_d$.  The integral mode includes the term
  $\int_0^t()dt$.}  
u(t) & = {\color{blue} K \left\{ e(t)+\frac
  1{\tau_i}\int_0^te(t)dt + \tau_d\frac{de(t)}{dt}
\right\}}
\end{align*}
\end{frame}

% simple animation
\begin{frame}
  \frametitle{From Spoken Tutorial on Letter
    Writing}
National Mission on Education through ICT, with an
outlay of Rs. 4,600 crore (\$ 1 billion): \pause
\begin{enumerate}
\item<+-|alert@+> Rs. 1,800 crore has been reserved for
  content generation and the rest to establish good
  connectivity in all 20,000 colleges and 200
  universities.
\item<+-|alert@+> Support for all good proposals,
  including those from private colleges.
\item<+-|alert@+> All products funded by this mission
  will be delivered as open source.
\item<+-|alert@+> Web based support through
  www.sakshat.ac.in.
\end{enumerate}
\end{frame}

% including figures
\begin{frame}
\frametitle{Example of a Figure}
\begin{center}
\includegraphics[width=\linewidth]{iitb}
\end{center}
\end{frame}

% guidelines for including figures
\begin{frame}
\frametitle{Hints for Including Figures}
\begin{itemize}
\item<+-> Do not use floated environments in
  presentations 
  - \\ {\tt begin\{figure\}}, {\tt end\{figure\}}, 
  etc. 
\item<+-> Use {\tt includegraphics} directly
\item<+-> Do not include {\tt caption}, figure
  number, etc.
\item<+-> The audience will not remember figure
  numbers, any way
\item<+-> If you want to refer to a previously
  shown figure, show it again - do not refer to it
  by number
\end{itemize}
\end{frame}

% two columns
\begin{frame}
\frametitle{Two Columns}
\begin{minipage}[c]{0.45\textwidth}
\begin{enumerate}
\item<+-|alert@+> Rs. 1,800 crore has been reserved for
  content generation and the rest to establish good
  connectivity in all 20,000 colleges and 200
  universities.
\item<+-|alert@+> Support for all good proposals,
  including those from private colleges.
\end{enumerate}
\end{minipage} \pause
\hfill
\begin{minipage}[c]{0.45\textwidth}
\includegraphics[width=\linewidth]{iitb}
\end{minipage}
\end{frame}

% example of a Table
\begin{frame}
\frametitle{Table: From Spoken Tutorial on Tables} 
\begin{center}
\begin{tabular}{||l|c|c|c|r|}\hline
\multicolumn 2 {||c|}{Fruit details} & 
\multicolumn 3 {c|}{Cost calculations} \\ \hline
Fruit & Type & No. of units & cost/unit & cost (Rs.) \\
\hline 
Mango & Malgoa & 18 & 50 &  \\ \cline{2-4}
      & Alfonso & 2 & 300 & 1,500 \\ \hline
Jackfruit & Kolli Hills & 10 & 50 & 500 \\ \hline
Banana & Green & 10 & 20 & 200 \\ \hline
\multicolumn 4{||r|}{Total cost (Rs.)} & 2,200 \\
\hline 
\end{tabular}
\end{center}
\end{frame}

% guidelines for including Tables
\begin{frame}
\frametitle{Hints for Including Tables}
\begin{itemize}
\item<+-> Do not use floated environments in
  presentations 
  - \\ {\tt begin\{table\}}, {\tt end\{table\}}, 
  etc. 
\item<+-> Use {\tt tabular} directly
\item<+-> Do not include {\tt caption}, Table
  number, etc.
\item<+-> The audience will not remember Table
  numbers, any way
\end{itemize}
\end{frame}

% using fragile with verbatim
\begin{frame}[fragile]
\frametitle{Example of Verbatim to Illustrate Scilab}
\begin{verbatim}
-->a = 1:5, b = 1:2:9
\end{verbatim}
{\color{blue}
\begin{verbatim}
 a  =
 !   1.    2.    3.    4.    5. !
 b  =
 !   1.    3.    5.    7.    9. !
\end{verbatim}}

\begin{verbatim}
-->a - 2
\end{verbatim}
{\color{blue}
\begin{verbatim}
ans  =
 ! - 1.    0.    1.    2.    3. !
\end{verbatim} }
\end{frame}

% where to get more information
\begin{frame}
\frametitle{Where to Get more Information on
  Beamer?} 
\begin{itemize}
\item<+-|alert@+> Authoritative source: user guide to the Beamer class
\item<+-|alert@+> Author: Till Tantau
\item<+-|alert@+> File name: 
  {\color{magenta} \tt beameruserguide.pdf} 
\item<+-|alert@+> email: tantau@users.sourceforge.net 
\item<+-|alert@+> At the time of creating this
  tutorial, available at \\
{\color{magenta} \small http://www.ctan.org/tex-archive/macros/latex/contrib/beamer/doc/beameruserguide.pdf}
\item<+-|alert@+> The web page for the Beamer
  project is
{\color{magenta} http://latex-beamer.sourceforge.net}
\end{itemize}
\end{frame}

\begin{frame} 
\frametitle{There Is No Largest Prime Number} 
\framesubtitle{The proof uses \textit{reductio ad absurdum}.} 
\begin{theorem} 
There is no largest prime number. 
\end{theorem} 
\begin{proof} 
\begin{enumerate} 
\item<1-| alert@1> Suppose $p$ were the largest prime number. 
\item<2-> Let $q$ be the product of the first $p$ numbers. 
\item<3-> Then $q+1$ is not divisible by any of them. 
\item<1-> Thus $q+1$ is also prime and greater than $p$.\qedhere 
\end{enumerate} 
\end{proof} 
\end{frame} 

% acknowledgement
\begin{frame}
\frametitle{Acknowledgement}
\begin{itemize}
\item Funding to create this tutorial has came from
  the \\ National Mission on Education through ICT
\item Web site for this mission:
  {\color{magenta} http://sakshat.ac.in}
\end{itemize}
\end{frame}

% feedback
\begin{frame}
\frametitle{Feedback}
\begin{itemize}
\item We have come to end of this tutorial
\item Thanks for joining 
\item Please give your feedback at kannan@iitb.ac.in
\end{itemize}
\end{frame}

\end{document}

